\documentclass[10pt,a4paper]{amsart}
\usepackage{amsmath,amssymb,amsthm,graphicx}
% Retained minimal preamble for the Stacks Project chapter more-algebra.tex (More on Algebra).
% The chapter file pinned-original-source.tex carries only \input{preamble} on its first line, so no
% preamble text lives inside the chapter itself.  The Stacks Project uses one single shared preamble for
% every chapter; the snapshot stages it identically as sources/stacks-chow/preamble.tex and
% sources/stacks-etale-cohomology/preamble.tex.  The block below is copied mechanically, line for line,
% from sources/stacks-chow/preamble.tex lines 14-18 (the Xy-pic setup) and 154-217 (the native theorem
% declarations together with the mathematical macros).  Nothing was added, removed or redefined; the
% wrapper class and the amsmath/amssymb/amsthm/graphicx base packages are supplied by the assembler.
% Chapter: more-algebra.tex -- More on Algebra.
% For commutative diagrams we use Xy-pic
\usepackage[all]{xy}

% We use 2cell for 2-commutative diagrams.
\xyoption{2cell}

% Theorem environments.
%
\theoremstyle{plain}
\newtheorem{theorem}[subsection]{Theorem}
\newtheorem{proposition}[subsection]{Proposition}
\newtheorem{lemma}[subsection]{Lemma}

\theoremstyle{definition}
\newtheorem{definition}[subsection]{Definition}
\newtheorem{example}[subsection]{Example}
\newtheorem{exercise}[subsection]{Exercise}
\newtheorem{situation}[subsection]{Situation}

\theoremstyle{remark}
\newtheorem{remark}[subsection]{Remark}
\newtheorem{remarks}[subsection]{Remarks}

\numberwithin{equation}{subsection}

% Macros
%
\def\lim{\mathop{\mathrm{lim}}\nolimits}
\def\colim{\mathop{\mathrm{colim}}\nolimits}
\def\Spec{\mathop{\mathrm{Spec}}}
\def\Hom{\mathop{\mathrm{Hom}}\nolimits}
\def\Ext{\mathop{\mathrm{Ext}}\nolimits}
\def\SheafHom{\mathop{\mathcal{H}\!\mathit{om}}\nolimits}
\def\SheafExt{\mathop{\mathcal{E}\!\mathit{xt}}\nolimits}
\def\Sch{\mathit{Sch}}
\def\Mor{\mathop{\mathrm{Mor}}\nolimits}
\def\Ob{\mathop{\mathrm{Ob}}\nolimits}
\def\Sh{\mathop{\mathit{Sh}}\nolimits}
\def\NL{\mathop{N\!L}\nolimits}
\def\CH{\mathop{\mathrm{CH}}\nolimits}
\def\proetale{{pro\text{-}\acute{e}tale}}
\def\etale{{\acute{e}tale}}
\def\QCoh{\mathit{QCoh}}
\def\Ker{\mathop{\mathrm{Ker}}}
\def\Im{\mathop{\mathrm{Im}}}
\def\Coker{\mathop{\mathrm{Coker}}}
\def\Coim{\mathop{\mathrm{Coim}}}

% Boxtimes
%
\DeclareMathSymbol{\boxtimes}{\mathbin}{AMSa}{"02}

%
% Macros for moduli stacks/spaces
%
\def\QCohstack{\mathcal{QC}\!\mathit{oh}}
\def\Cohstack{\mathcal{C}\!\mathit{oh}}
\def\Spacesstack{\mathcal{S}\!\mathit{paces}}
\def\Quotfunctor{\mathrm{Quot}}
\def\Hilbfunctor{\mathrm{Hilb}}
\def\Curvesstack{\mathcal{C}\!\mathit{urves}}
\def\Polarizedstack{\mathcal{P}\!\mathit{olarized}}
\def\Complexesstack{\mathcal{C}\!\mathit{omplexes}}
% \Pic is the operator that assigns to X its picard group, usage \Pic(X)
% \Picardstack_{X/B} denotes the Picard stack of X over B
% \Picardfunctor_{X/B} denotes the Picard functor of X over B
\def\Pic{\mathop{\mathrm{Pic}}\nolimits}
\def\Picardstack{\mathcal{P}\!\mathit{ic}}
\def\Picardfunctor{\mathrm{Pic}}
\def\Deformationcategory{\mathcal{D}\!\mathit{ef}}

% --- appended 2026-10-12 (Chief): the project own definitions of the reference/slogan/history
% display environments, copied mechanically from the shared Stacks preamble lines 8-12
% (\usepackage{verbatim} plus the three \newenvironment lines, which the chapter preamble keeps
% verbatim from the project preamble). Nothing added, removed or redefined: these environments
% carry no mathematical content and the project itself typesets them as comments. Without these
% three lines a counted statement whose own text uses them fails with Environment slogan undefined.
\usepackage{verbatim}
\newenvironment{reference}{\comment}{\endcomment}
\newenvironment{slogan}{\comment}{\endcomment}
\newenvironment{history}{\comment}{\endcomment}
\usepackage{cleveref}
\crefname{originallemma}{original lemma}{original lemmas}
\Crefname{originallemma}{Original lemma}{Original lemmas}
\makeatletter
\newlabel{algebra-lemma-etale-over-field}{{00U3}{source}{}{sourcepointer-1}{}}
\newlabel{algebra-lemma-etale-over-field@cref}{{[originallemma][0][]00U3}{[1][0][]source}{}{}{}}
\makeatother
\crefname{originallemma}{original lemma}{original lemmas}
\Crefname{originallemma}{Original lemma}{Original lemmas}
\makeatletter
\newlabel{algebra-lemma-finite-residue-extension-closed}{{00GA}{source}{}{sourcepointer-2}{}}
\newlabel{algebra-lemma-finite-residue-extension-closed@cref}{{[originallemma][0][]00GA}{[1][0][]source}{}{}{}}
\makeatother
\crefname{originallemma}{original lemma}{original lemmas}
\Crefname{originallemma}{Original lemma}{Original lemmas}
\makeatletter
\newlabel{algebra-lemma-injective-minimal-primes-in-image}{{00FK}{source}{}{sourcepointer-3}{}}
\newlabel{algebra-lemma-injective-minimal-primes-in-image@cref}{{[originallemma][0][]00FK}{[1][0][]source}{}{}{}}
\makeatother
\crefname{originallemma}{original lemma}{original lemmas}
\Crefname{originallemma}{Original lemma}{Original lemmas}
\makeatletter
\newlabel{algebra-proposition-dimension-zero-ring}{{00KJ}{source}{}{sourcepointer-4}{}}
\newlabel{algebra-proposition-dimension-zero-ring@cref}{{[originallemma][0][]00KJ}{[1][0][]source}{}{}{}}
\makeatother
\crefname{originallemma}{original lemma}{original lemmas}
\Crefname{originallemma}{Original lemma}{Original lemmas}
\makeatletter
\newlabel{lemma-absolutely-flat}{{30250}{source}{}{sourcepointer-5}{}}
\newlabel{lemma-absolutely-flat@cref}{{[originallemma][0][]30250}{[1][0][]source}{}{}{}}
\makeatother
\crefname{originallemma}{original lemma}{original lemmas}
\Crefname{originallemma}{Original lemma}{Original lemmas}
\makeatletter
\newlabel{lemma-absolutely-flat-fields}{{30488}{source}{}{sourcepointer-6}{}}
\newlabel{lemma-absolutely-flat-fields@cref}{{[originallemma][0][]30488}{[1][0][]source}{}{}{}}
\makeatother
\crefname{originallemma}{original lemma}{original lemmas}
\Crefname{originallemma}{Original lemma}{Original lemmas}
\makeatletter
\newlabel{lemma-absolutely-flat-over-absolutely-flat}{{30323}{source}{}{sourcepointer-7}{}}
\newlabel{lemma-absolutely-flat-over-absolutely-flat@cref}{{[originallemma][0][]30323}{[1][0][]source}{}{}{}}
\makeatother
\crefname{originallemma}{original lemma}{original lemmas}
\Crefname{originallemma}{Original lemma}{Original lemmas}
\makeatletter
\newlabel{lemma-composition-weakly-etale}{{30347}{source}{}{sourcepointer-8}{}}
\newlabel{lemma-composition-weakly-etale@cref}{{[originallemma][0][]30347}{[1][0][]source}{}{}{}}
\makeatother
\crefname{originallemma}{original lemma}{original lemmas}
\Crefname{originallemma}{Original lemma}{Original lemmas}
\makeatletter
\newlabel{lemma-when-weakly-etale}{{30454}{source}{}{sourcepointer-9}{}}
\newlabel{lemma-when-weakly-etale@cref}{{[originallemma][0][]30454}{[1][0][]source}{}{}{}}
\makeatother
\title{A tensor-square flat algebra over a field is weakly etale and a filtered colimit of etale algebras}
\author{The Stacks Project Authors}
\date{stacks-project:more-algebra.tex:local-snapshot-20261012; GNU FDL 1.2 or later}
\begin{document}
\maketitle
\setcounter{subsection}{1}
\noindent\textit{Editorial note.} All delivered mathematical text is the author's own, copied from the pinned Stacks Project chapter file more-algebra.tex; no statement or proof step was rewritten, reordered or removed.\par\medskip
\noindent\textit{Editorial note.} The chapter file pinned-original-source.tex carries only \textbackslash{}input\{preamble\} on its first line, so no preamble text lives inside the chapter. The retained file sources/stacks-more-algebra/retained-preamble.tex is a mechanical, line-for-line copy of the project's single shared preamble as staged in the snapshot; nothing was added, removed or redefined and no package was added.\par\medskip
\noindent\textit{Editorial note.} The source URL is the source-level Stacks Project URL, as required when more than one card is produced from one source; the per-claim tag is carried by claim\_key.\par\medskip
\noindent\textit{Editorial note.} The author writes 'hence B is a absolutely flat ring' and 'a finitely generated K-sub algebra'; both phrasings are reproduced unchanged. The nine references inside the counted proof span are recorded as non-load-bearing navigation pointers.\par\medskip
\noindent\textit{Allowed standard prerequisites (editorial).}\begin{itemize}
\item Absolutely flat rings: reduced, with every local ring a field. invoked at proof lines 30522-30525. Gives reducedness of B and the field property of its local rings, which the rest of the argument uses.
\item Weakly etale maps and their composition, and weakly etale fields as separable algebraic extensions. used at proof lines 30528-30532. Converts the hypothesis into the separability of the residue field extensions.
\item Dimension-zero rings as products of their local rings, and etale algebras over a field as reduced rings with finite separable residue fields. used at proof lines 30546-30552. Is the criterion that turns the dimension-zero conclusion into etaleness of A over K.
\end{itemize}\medskip
\ifdefined\hypertarget\hypertarget{sourcepointer-1}{}\fi
\noindent\textit{Source reference.} Original lemma 00U3: focused-seeds/AGP/sources/stacks-snapshot-20261011/tags.txt, source lines 1097--1097. Cross-chapter navigation pointer into the Stacks Project; the standard fact it supplies is applied inside the counted proof span and the argument resting on it stays there.\par\medskip
\ifdefined\hypertarget\hypertarget{sourcepointer-2}{}\fi
\noindent\textit{Source reference.} Original lemma 00GA: focused-seeds/AGP/sources/stacks-snapshot-20261011/tags.txt, source lines 600--600. Cross-chapter navigation pointer into the Stacks Project; the standard fact it supplies is applied inside the counted proof span and the argument resting on it stays there.\par\medskip
\ifdefined\hypertarget\hypertarget{sourcepointer-3}{}\fi
\noindent\textit{Source reference.} Original lemma 00FK: focused-seeds/AGP/sources/stacks-snapshot-20261011/tags.txt, source lines 574--574. Cross-chapter navigation pointer into the Stacks Project; the standard fact it supplies is applied inside the counted proof span and the argument resting on it stays there.\par\medskip
\ifdefined\hypertarget\hypertarget{sourcepointer-4}{}\fi
\noindent\textit{Source reference.} Original lemma 00KJ: focused-seeds/AGP/sources/stacks-snapshot-20261011/tags.txt, source lines 753--753. Cross-chapter navigation pointer into the Stacks Project; the standard fact it supplies is applied inside the counted proof span and the argument resting on it stays there.\par\medskip
\ifdefined\hypertarget\hypertarget{sourcepointer-5}{}\fi
\noindent\textit{Source reference.} Original lemma 30250: sources/stacks-more-algebra/pinned-original-source.tex, source lines 30250--30250. Same-chapter navigation pointer; the statement it names is standard for the section and the argument that consumes it stays inside the counted proof span.\par\medskip
\ifdefined\hypertarget\hypertarget{sourcepointer-6}{}\fi
\noindent\textit{Source reference.} Original lemma 30488: sources/stacks-more-algebra/pinned-original-source.tex, source lines 30488--30488. Same-chapter navigation pointer; the statement it names is standard for the section and the argument that consumes it stays inside the counted proof span.\par\medskip
\ifdefined\hypertarget\hypertarget{sourcepointer-7}{}\fi
\noindent\textit{Source reference.} Original lemma 30323: sources/stacks-more-algebra/pinned-original-source.tex, source lines 30323--30323. Same-chapter navigation pointer; the statement it names is standard for the section and the argument that consumes it stays inside the counted proof span.\par\medskip
\ifdefined\hypertarget\hypertarget{sourcepointer-8}{}\fi
\noindent\textit{Source reference.} Original lemma 30347: sources/stacks-more-algebra/pinned-original-source.tex, source lines 30347--30347. Same-chapter navigation pointer; the statement it names is standard for the section and the argument that consumes it stays inside the counted proof span.\par\medskip
\ifdefined\hypertarget\hypertarget{sourcepointer-9}{}\fi
\noindent\textit{Source reference.} Original lemma 30454: sources/stacks-more-algebra/pinned-original-source.tex, source lines 30454--30454. Same-chapter navigation pointer; the statement it names is standard for the section and the argument that consumes it stays inside the counted proof span.\par\medskip
\begin{lemma}
\label{lemma-absolutely-flat-over-field}
Let $B$ be an algebra over a field $K$. The following are
equivalent
\begin{enumerate}
\item $B \otimes_K B \to B$ is flat,
\item $K \to B$ is weakly \'etale, and
\item $B$ is a filtered colimit of \'etale $K$-algebras.
\end{enumerate}
Moreover, every finitely generated $K$-subalgebra of $B$
is \'etale over $K$.
\end{lemma}

\begin{proof}
Parts (1) and (2) are equivalent because every $K$-algebra is flat over $K$.
Part (3) implies (1) and (2) by Lemma \ref{lemma-when-weakly-etale}

\medskip\noindent
Assume (1) and (2) hold. We will prove (3) and the finite statement of
the lemma. A field is absolutely flat ring, hence $B$ is a absolutely
flat ring by Lemma \ref{lemma-absolutely-flat-over-absolutely-flat}.
Hence $B$ is reduced and every local
ring is a field, see Lemma \ref{lemma-absolutely-flat}.

\medskip\noindent
Let $\mathfrak q \subset B$ be a prime. The ring map
$B \to B_\mathfrak q$ is weakly \'etale, hence $B_\mathfrak q$
is weakly \'etale over $K$ (Lemma \ref{lemma-composition-weakly-etale}).
Thus $B_\mathfrak q$ is a separable algebraic extension of $K$ by
Lemma \ref{lemma-absolutely-flat-fields}.

\medskip\noindent
Let $K \subset A \subset B$ be a finitely generated $K$-sub algebra.
We will show that $A$ is \'etale over $K$ which will finish the proof
of the lemma.
Then every minimal prime $\mathfrak p \subset A$ is the image of a prime
$\mathfrak q$ of $B$, see
Algebra, Lemma \ref{algebra-lemma-injective-minimal-primes-in-image}.
Thus $\kappa(\mathfrak p)$ as a subfield of
$B_\mathfrak q = \kappa(\mathfrak q)$ is separable algebraic over $K$.
Hence every generic point of $\Spec(A)$
is closed (Algebra, Lemma \ref{algebra-lemma-finite-residue-extension-closed}).
Thus $\dim(A) = 0$.
Then $A$ is the product of its local rings, e.g., by
Algebra, Proposition \ref{algebra-proposition-dimension-zero-ring}.
Moreover, since $A$ is reduced, all local rings are equal
to their residue fields which are finite separable over $K$.
This means that $A$ is \'etale over $K$ by
Algebra, Lemma \ref{algebra-lemma-etale-over-field}
and finishes the proof.
\end{proof}

\end{document}
